\chapter{总结与展望}
\section {研究总结}
具身视觉导航作为智能机器人感知与决策的核心技术，是实现家庭服务、工业巡检、公共服务等场景自主化的关键基础。
随着应用场景的不断拓展，导航任务对场景建模的精度、目标识别的鲁棒性与导航决策的效率提出了更高要求。
其中，实例图像目标导航（IIN）要求智能体仅通过单张目标物体图像，在未知环境中定位并导航至该特定实例，需解决位置环境探索、同类别实例区分抗干扰三大核心挑战。
然而，现有主流方法普遍采用二维鸟瞰图（BEV）作为地图表示，虽具备内存占用低、路径规划便捷的优势，但存在三维几何信息缺失、实例纹理特征丢失等固有缺陷，难以支撑高精度的实例级目标定位与匹配，成为制约 IIN 任务性能提升的核心瓶颈。
3D 高斯泼溅（3DGS）技术凭借其显式的三维几何表示、高保真的纹理渲染能力与实时可微分渲染特性，为突破上述瓶颈提供了全新的技术路径。
GaussNav\cite{gaussnav}首先使用了 3D 高斯地图表示，将目标物体的语义信息与几何信息联合表示，避免了每回合都需要重新构建地图的计算成本，提升了目标定位效率和准确度。
但语义信息仅参与高斯标签的赋予阶段，未参与后续的监督优化阶段，导致一定程度上的语义覆盖不完整、语义信息缺失等问题。
针对上述问题，本文基于 GaussNav\cite{gaussnav} 基础框架，融合 GSplatLoc\cite{gsplatloc} 的位姿精修算法，构建了一套集几何、语义与实例特征于一体的语义高斯导航系统，实现了仅输入单张目标图像即可完成目标实例定位与高效导航的完整流程。
研究从地图表示设计、在线构建机制、目标定位算法与系统性能验证四个层面系统开展，取得了以下主要成果与创新：
\begin {enumerate}
\item \textbf {语义增强的高斯地图表示与联合优化机制}：针对传统 BEV 地图无法保留物体三维几何与纹理细节的问题，本文在 GaussNav\cite{gaussnav} 框架基础上新增语义渲染通道，采用 YOLO26\cite{Yolo26} 语义分割网络输出的像素级语义信息，在高斯初始化与致密化的全流程中为每个高斯赋予多维语义向量，并将语义损失与 RGB、深度损失共同纳入高斯参数的监督优化过程。该机制实现了几何、纹理与语义信息的端到端联合学习，显著提升了语义高斯地图的实例区分能力与跨视角匹配精度。
% \item \textbf {增量式在线高斯构建与并行化处理}：针对现有高斯地图构建需离线批量处理观测数据、无法支持实时导航的问题，本文设计了高并行的在线高斯构建流水线。智能体在探索过程中采集 RGB-D 图像的同时，实时完成语义分割、高斯致密化与参数更新，实现了 “边探索、边建图、边可用” 的闭环流程。该策略简化了任务执行的先后依赖，大幅缩短了从环境探索到导航可用的时间延迟，提升了系统的整体运行效率。
\item \textbf {两阶段目标位姿精修与高效匹配策略}：针对传统方法目标定位精度不足、需多次验证导致导航效率低下的问题，本文融合 GSplatLoc\cite{gsplatloc} 的可微分渲染位姿精修算法，提出 “初始粗定位 + 渲染精修” 的两阶段目标定位框架。首先通过多视角渲染与关键点匹配实现目标实例的粗定位，再利用可微分渲染优化目标相机位姿，大幅提升了目标定位的精度。同时，位姿精修过程仅需单次前向渲染即可完成梯度计算与参数更新，有效降低了内存开销与计算延迟。
% \item \textbf {系统性实验验证与性能对比分析}：基于 Habitat\cite{habitat-sim} 高保真仿真平台与 HM3D\cite{HM3D} 大规模真实场景数据集，本文设计了涵盖不同场景复杂度、不同目标类别与不同初始状态的系统性实验，从成功率、SPL、运行效率与内存占用四个维度全面评估了所提方法的性能。实验结果表明，本文方法在 HM3D\cite{HM3D} 数据集上的 SPL 指标从原有 GaussNav 的 0.578 进一步提升至 0.612，同时保持了 20FPS 以上的实时运行帧率，在定位精度与导航效率上均显著优于现有 BEV 地图基的代表方法，验证了所提框架的有效性与优越性。
\item \textbf{目标导航：} 利用优化后的语义高斯表示，结合基于LightGlue+Disk\cite{lightglue,disk}的特征提取与匹配方案的目标位姿估计，实现了目标导航。
\end {enumerate}
通过以上研究，本文构建了基于 3D 高斯泼溅的实例级视觉导航完整技术体系，突破了传统 BEV 地图在三维信息保留与实例特征表征上的核心局限，实现了从 “语义级导航” 到 “高精度实例级导航” 的跨越。
研究不仅在理论上丰富了具身视觉导航的地图表示方法与目标定位技术，也在实践中为服务机器人的物品查找、室内巡检等实际应用提供了可行的技术方案，具备良好的产业化应用前景。

\section {研究展望}
% 然而，需要指出的是，尽管本文方法在仿真环境中已取得良好效果，但在实际应用中，真实场景的光照变化、纹理缺失、动态物体干扰以及传感器噪声等因素仍可能影响高斯地图构建质量与目标匹配精度。
% 同时，当前方法仍依赖全局预探索完成地图构建，无法直接应用于动态变化的未知场景，这些问题均有待在后续研究中进一步解决。

尽管本文围绕实例图像目标导航任务，提出了基于语义高斯的导航方法体系并在仿真环境中验证了其有效性，但在实际应用中，真实场景的光照变化、纹理缺失、动态物体干扰以及传感器噪声等因素仍可能影响高斯地图构建质量与目标匹配精度，仍存在诸多值得进一步深入探究的问题与挑战。
结合当前 3D 高斯泼溅技术与具身智能的发展趋势，本文认为后续研究可在以下几个方面继续深入，推动实例级视觉导航技术迈向更高水平。

首先，提升方法在真实物理环境中的适应性与泛化能力是未来研究的首要任务。
尽管本文方法已在 HM3D 仿真数据集上展现出优异的性能，但仿真环境与真实世界之间存在不可避免的域偏移，因此，未来研究可探索结合域自适应、元学习与在线增量学习技术，实现跨场景、跨设备的快速适配，增强方法在复杂真实环境中的可靠性。

其次，引入可微分物理与动态高斯建模机制，实现更精准的场景表示与位姿估计，也是未来值得重点关注的方向。
未来可将可微分物理引擎与 3D 高斯泼溅相结合，引入物理一致性约束优化高斯参数，提升场景重建的精度与合理性。
同时，研究动态 3D 高斯表示方法，实现对场景中运动物体的实时建模与跟踪，拓展导航系统在动态环境中的应用能力。

最后，实现大场景轻量化在线构建与人机协同交互是技术走向实际应用的必要环节。
未来研究可探索高斯压缩、稀疏化与层次化表示技术，降低大场景地图的内存开销；
研究增量式高斯构建与回环检测算法，实现边导航边建图的未知环境探索模式。

综上，基于 3D 高斯泼溅的实例级视觉导航仍处于快速发展阶段，面向真实复杂环境的鲁棒导航、动态场景建模与多任务协同等问题依然充满挑战与机遇。
未来，随着多模态感知、可微分物理、大模型具身智能等技术的不断进步，语义高斯导航技术有望实现从实验室研究向实际应用的全面转化，为智能服务机器人在家庭、医院、商场等场景的广泛应用提供坚实的技术支撑。
